\documentclass[10pt,twocolumn]{article}
\usepackage[T1]{fontenc}
\usepackage[utf8]{inputenc}
\usepackage{lmodern}
\usepackage[margin=1.8cm,top=2.4cm,bottom=2cm,columnsep=0.6cm]{geometry}
\usepackage{fancyhdr}
\usepackage{titlesec}
\usepackage{booktabs}
\usepackage{amsmath,amssymb,amsfonts}
\usepackage{microtype}
\usepackage{xcolor}
\usepackage{mdframed}
\usepackage{parskip}
\usepackage{enumitem}
\usepackage{hyperref}

\definecolor{rulered}{RGB}{180,30,30}
\definecolor{boxgray}{RGB}{245,245,245}
\definecolor{keywordgray}{RGB}{100,100,100}

\pagestyle{fancy}
\fancyhf{}
\fancyhead[L]{\textsc{\small Claude Mag}}
\fancyhead[C]{\small\textit{Machine Learning Research Digest}}
\fancyhead[R]{\small 2026-11-02}
\fancyfoot[C]{\small\thepage}
\renewcommand{\headrulewidth}{0.8pt}

\titleformat{\section}{\normalsize\bfseries\color{rulered}}{}{0pt}{}%
  [\vspace{-4pt}\color{rulered}\rule{\columnwidth}{0.4pt}]
\titlespacing{\section}{0pt}{8pt}{4pt}

\hypersetup{colorlinks=true,urlcolor=rulered,linkcolor=black}

\begin{document}
\setlength{\parindent}{0pt}
\setlength{\parskip}{4pt}

{\small\color{keywordgray}\textbf{Keywords:}
  vision-and-language navigation \textbullet{}
  hierarchical reinforcement learning \textbullet{}
  embodied AI \textbullet{}
  continuous environments}\par\vspace{4pt}

{\LARGE\bfseries\raggedright
  MATCH: Execution-Aware Hierarchical Reinforcement Learning
  for Vision-and-Language Navigation}\par\vspace{4pt}

{\large\itshape\raggedright
  Joint End-to-End Training of Planner and Executor Sets
  New State-of-the-Art on R2R-CE and RxR-CE at 0.5B Parameters}\par\vspace{6pt}

{\small\textbf{By} Shuhao Ye, Yufei Wei, Demao Wang,
  Shichao Zhai, Rong Xiong, Yue Wang}\par\vspace{2pt}

{\small\color{keywordgray}arXiv:2610.38248 \textbullet{} cs.RO
  \textbullet{} submitted 28 Oct 2026}
\par\vspace{2pt}\noindent{\color{rulered}\rule{\columnwidth}{0.6pt}}\vspace{6pt}

Hierarchical VLN agents that optimize the high-level planner while
keeping the low-level executor fixed suffer from execution blindness:
the planner proposes waypoints without knowing whether the executor
can reach them. MATCH closes this loop by jointly training a
language-conditioned planner and a goal-conditioned executor end-to-end
with reinforcement learning, routing execution feedback back to inform
future waypoint selection. With only 0.5B parameters, MATCH sets
new state-of-the-art on R2R-CE (62.4\% SR) and RxR-CE (55.8\% SR).

\begin{mdframed}[backgroundcolor=boxgray,linecolor=rulered,linewidth=1pt,
    innerleftmargin=6pt,innerrightmargin=6pt,
    innertopmargin=6pt,innerbottommargin=6pt]
\textbf{\small AT A GLANCE}\par\vspace{4pt}
\begin{description}[leftmargin=0pt,labelsep=4pt,itemsep=2pt,parsep=0pt,
    font=\small\bfseries]
  \item[Problem:] In VLN-CE, hierarchical planners propose waypoints
    without knowing what the executor can achieve, causing cascading
    failures.
  \item[Why it matters:] Embodied navigation in continuous 3D
    environments with language instructions is a core capability
    for household robots and assistive AI.
  \item[Method:] End-to-end RL training of planner and executor;
    execution-status feedback from executor to planner at each step;
    shared visual backbone.
  \item[Result:] 62.4\% SR on R2R-CE (val unseen) and 55.8\% SR on
    RxR-CE (val unseen), surpassing all prior methods.
  \item[Evidence:] R2R-CE and RxR-CE benchmarks, val unseen and
    test splits; ablations over execution feedback and joint training.
\end{description}
\end{mdframed}

\section{The Big Picture}

Vision-and-Language Navigation (VLN) tests an agent's ability to
follow free-form instructions like ``walk past the sofa, turn right at
the kitchen, and stop in front of the bookshelf'' in a photo-realistic
3D environment. In Continuous Environments (VLN-CE), the agent must
produce raw motor actions---translation velocity and rotation---rather
than selecting named waypoints on a graph, making the problem
significantly harder.

Hierarchical approaches decompose VLN-CE into two sub-problems:
a \emph{planner} that parses the instruction and generates a
sequence of subgoals (waypoints), and an \emph{executor} that
drives the robot from its current position to each waypoint.
This decomposition is intuitive but introduces a structural mismatch:
the planner has no direct signal about whether its proposed waypoints
are achievable given the current room geometry and the executor's
capabilities.

\section{Why Existing Approaches Fall Short}

\textbf{Fixed executor, trained planner:}
Most prior hierarchical VLN-CE systems pretrain or fix the low-level
executor (typically a model-predictive controller or a simple
goal-conditioned policy) and focus all learning effort on the
high-level planner. The planner therefore cannot receive gradient
signal from execution failures and learns only from its own
prediction errors.

\textbf{Mismatched subgoal abstraction:}
Even when both components are trained, they are typically trained
separately on different objectives (language-goal matching for the
planner; reaching accuracy for the executor), meaning there is no
shared objective that penalizes the planner for proposing unreachable
waypoints.

\textbf{No execution feedback:}
The planner selects its next waypoint based on the instruction and
the current visual observation, but not on how the last waypoint
was executed. Information about partial failures (executor reached
70\% of the waypoint but collided) is discarded between planning
steps.

\section{Core Mechanism}

\textbf{Language-Conditioned Planner.} Given the full instruction
$I$ and the current visual observation $o_t$, the planner
$\pi_P(w_t \mid I, o_t, h_t^P)$ produces a waypoint $w_t$ as a
displacement vector relative to the current agent pose, where
$h_t^P$ is the planner's recurrent state.

\textbf{Goal-Conditioned Executor.} Given the current observation
$o_t^{\text{low}}$ and the target waypoint $w_t$, the executor
$\pi_E(a_t \mid o_t^{\text{low}}, w_t, h_t^E)$ produces a
primitive action $a_t$ (linear velocity, angular velocity).

\textbf{Execution-Status Feedback.} After the executor attempts
waypoint $w_t$, the achieved displacement and an obstacle-collision
indicator are compressed into an execution-status vector $s_t$ and
appended to the planner's observation at step $t+1$:
$h_{t+1}^P = \text{RNN}(h_t^P, [o_{t+1}, s_t])$.

\textbf{Joint RL Training.} The planner and executor share a single
reward: the navigation success reward at the episode terminus,
discounted back through the waypoint sequence. Both are updated by
PPO in a shared replay buffer.

\section{Technical Formulation}

The planner's policy is:
\[
w_t = \pi_P\!\left(I,\, o_t,\, h_t^P,\, s_{t-1}\right)
\]
The executor's policy unrolls for $T_{\text{exec}}$ steps per waypoint:
\[
a_\tau = \pi_E\!\left(o_\tau^{\text{low}},\, w_t - p_\tau,\, h_\tau^E\right)
\]
where $p_\tau$ is the agent's current position.

The reward shaping for the planner uses a distance-based dense reward:
\[
r_t^P = d(p_t, g) - d(p_{t+1}, g) - \lambda_c \cdot \mathbf{1}[\text{collision}_t]
\]
where $g$ is the goal location and $\lambda_c$ penalizes collisions.
The executor receives a waypoint-reaching reward:
\[
r_\tau^E = \max\!\left(0,\; d(p_\tau, w_t) - d(p_{\tau+1}, w_t)\right)
\]

The full PPO objective optimizes $\pi_P$ and $\pi_E$ jointly:
\[
\mathcal{L} = \mathcal{L}_{\text{PPO}}(\pi_P) + \beta\,
  \mathcal{L}_{\text{PPO}}(\pi_E)
\]

\section{Learning or Inference Procedure}

\textbf{Pretraining.} The shared visual backbone (ViT-B) is initialized
from a pretrained visual-semantic embedding. The instruction encoder
is initialized from a compact BERT variant.

\textbf{Stage 1 -- Imitation Learning.} Both planner and executor are
pretrained with behavior cloning on expert trajectories for 10k steps.

\textbf{Stage 2 -- Joint RL Fine-tuning.} PPO is applied for 50k
environment steps. The execution-feedback vector is activated after
Stage 1, ensuring the planner does not receive noisy feedback
from an untrained executor.

\textbf{Inference.} The planner generates one waypoint at a time,
the executor runs for up to $T_{\text{exec}} = 20$ steps, and the
planner then receives the execution-status feedback and produces the
next waypoint.

\section{What the Guarantee Says}

No theoretical guarantee is claimed. The paper's argument is
empirical: execution-aware training produces better navigation
policies, and the ablation studies isolate which components
contribute. The success-rate metric on val unseen splits, which
feature environments not seen during training, provides evidence
of generalization.

\section{Experimental Findings}

\begin{itemize}[itemsep=2pt,parsep=0pt]
  \item \textbf{R2R-CE val unseen:} MATCH achieves
    \textbf{SR = 62.4\%}, \textbf{GP = 65.1\%},
    \textbf{NE = 3.8 m}, surpassing the prior SOTA
    (SR = 58.1\%) by 4.3 points.
  \item \textbf{RxR-CE val unseen:} MATCH achieves
    \textbf{SR = 55.8\%}, improving over the prior best (51.3\%)
    by 4.5 points.
  \item \textbf{Parameter efficiency:} MATCH achieves these results
    with \textbf{0.5B parameters}, outperforming prior
    methods that use 2--7B parameters.
  \item \textbf{Collision rate:} Joint RL training with the
    collision penalty reduces the executor's collision rate by
    31\% relative to the imitation-only baseline.
\end{itemize}

\section{Ablations and Interpretation}

\textbf{Execution feedback removed:}
Removing the execution-status vector and training jointly on reward
alone reduces SR on R2R-CE by 3.1 points, confirming that the
information channel---not just joint training---is critical.

\textbf{Fixed executor (standard hierarchical baseline):}
Fixing the executor at its imitation-learning checkpoint and training
only the planner reduces SR by 5.7 points, establishing the cost
of decoupled training.

\textbf{Imitation only (no RL):}
Without RL, SR is 54.6\%, 7.8 points below MATCH. RL is necessary
to generalize to execution scenarios not covered by expert data.

\textbf{Waypoint granularity:}
Decreasing the maximum number of waypoints per episode degrades
SR smoothly, suggesting the planner learns to use waypoints
efficiently rather than uniformly spacing them.

\section*{Reference}

Shuhao Ye, Yufei Wei, Demao Wang, Shichao Zhai, Rong Xiong, Yue Wang.
\textbf{MATCH: Execution-Aware Hierarchical Reinforcement Learning for
Vision-and-Language Navigation in Continuous Environment.}
arXiv:2610.38248. 2026.
\url{https://arxiv.org/abs/2610.38248}

\end{document}
